

\LabHeader{1}{ 1}

\begin{enumerate}
%\item Ввести с клавиатуры целое число (переменная \verb|n1|), вещественное число (переменная \verb|n2|) и слово (переменная \verb|w|).
%\begin{enumerate}
%\item Вывести значения переменных в одной строке, разделяя пробелом.
%\item Вывести значения переменных одно под другим.
%\item Вывести на экран значения переменных, сопровождая вывод комментариями, например:
%\begin{verbatim}
%Вы ввели три значения.
%Первое из них - ...
%Второе - ...
%А третье - ...
%\end{verbatim}
%\end{enumerate}
\item \textbf{День народження.} Складіть програму, яка запитує користувача дату народження

(Приклад: \verb|Коли у Вас День народження?|),

а потім після введення даних друкує на екрані повідомлення

(Приклад: \verb|Мій день народження ...|).

\item \textbf{Лінолеум.} Є шматок лінолеуму з розмірами $n \times m$. Скільки квадратів із стороною $k$ можна вирізати із нього?

{\itshape Намагайтеся супроводжувати введення та виведення текстовими повідомленнями.}

\item \textbf{Сума цифр.} Дано тризначне натуральне число. Знайти суму його цифр.

\item \textbf{Формула.} Дано деякі дійсні числа $x$ і $y$. Обчислити значення виразу$$
z=\dfrac{\sqrt{\mathstrut 1+\ln \left(1+x^2\right)}+\sqrt{\mathstrut 1+\dfrac{1}{1+|y|}}}{2x\,(1+y^2)}
$$
з точністю до сотих. Перевірка даних на коректність не потрібна.
\end{enumerate}



\cleardoublepage

\LabHeader{1}{ 2}

\begin{enumerate}
%\item Ввести с клавиатуры целое число (переменная \verb|n1|), вещественное число (переменная \verb|n2|) и слово (переменная \verb|w|).
%\begin{enumerate}
%\item Вывести значения переменных в одной строке, разделяя пробелом.
%\item Вывести значения переменных одно под другим.
%\item Вывести на экран значения переменных, сопровождая вывод комментариями, например:
%\begin{verbatim}
%Вы ввели три значения.
%Первое из них - ...
%Второе - ...
%А третье - ...
%\end{verbatim}
%\end{enumerate}

\item \textbf{Улюблена музика.} Складіть програму, яка запитує у користувача його улюблений музичний напрямок

(Приклад: \verb|Ваша улюблена музика:|),

а потім після введення даних друкує на екрані повідомлення

(Приклад: \ verb | Я люблю ... |).

\item \textbf{Куб.} Дана довжина ребра куба. Знайти об'єм куба та площу його бічної поверхні.

{\itshape Намагайтеся супроводжувати введення та виведення текстовими повідомленнями.}

\item \textbf{Маніпуляції з числом.} Дано тризначне натуральне число. У ньому закреслили першу цифру і приписали її наприкінці. Вивести на екран отримане число.

\item \textbf{Формула.} Дано деякі дійсні числа $x$ і $y$. Обчислити значення виразу
$$
z=\ln \left(2-\sin x^2\right)\cdot\dfrac{x\,(1+x)}{|y|+\dfrac{1}{\sqrt{\mathstrut \mathstrut x^2+10}}}
$$
з точністю до сотих. Перевірка даних на коректність не потрібна.
\end{enumerate}

\cleardoublepage

\LabHeader{1}{ 3}

\begin{enumerate}
%\item Ввести с клавиатуры целое число (переменная \verb|n1|), вещественное число (переменная \verb|n2|) и слово (переменная \verb|w|).
%\begin{enumerate}
%\item Вывести значения переменных в одной строке, разделяя пробелом.
%\item Вывести значения переменных одно под другим.
%\item Вывести на экран значения переменных, сопровождая вывод комментариями, например:
%\begin{verbatim}
%Вы ввели три значения.
%Первое из них - ...
%Второе - ...
%А третье - ...
%\end{verbatim}
%\end{enumerate}

\item \textbf{Улюблена книга.} Складіть програму, яка запитує у користувача назву його улюбленої книги

(Приклад: \verb|Ваша улюблена книга:|),

а потім після введення даних друкує на екрані повідомлення

(Приклад: \verb|Моя улюблена книга - ...|).

\item \textbf{Площа кільця.} Знайти площу кільця по заданим зовнішнім та внутрішнім радіусам.

{\itshape Намагайтеся супроводжувати введення та виведення текстовими повідомленнями.}

\item \textbf{Сума двох цифр.} Дано чотиризначне натуральне число. Знайти суму двох його перших цифр.

\item \textbf{Формула.} Дані деякі дійсні числа $a$ та $b$. Обчислити значення виразу
$$
z=\sin(a^2+b^2)\cdot\dfrac{\dfrac{a}{a^2+25}+b}{\sqrt{\mathstrut b}+\dfrac{a+b}{ 2}}
$$
з точністю до сотих. Перевірка даних на коректність не потрібна.
\end{enumerate}

\cleardoublepage

\LabHeader{1}{ 4}

\begin{enumerate}
%\item Ввести с клавиатуры целое число (переменная \verb|n1|), вещественное число (переменная \verb|n2|) и слово (переменная \verb|w|).
%\begin{enumerate}
%\item Вывести значения переменных в одной строке, разделяя пробелом.
%\item Вывести значения переменных одно под другим.
%\item Вывести на экран значения переменных, сопровождая вывод комментариями, например:
%\begin{verbatim}
%Вы ввели три значения.
%Первое из них - ...
%Второе - ...
%А третье - ...
%\end{verbatim}
%\end{enumerate}

\item \textbf{Привітання.} Складіть програму, яка запитує у користувача його ім'я та рік народження, а потім виводить повідомлення

\verb|Здрастуйте, ...! Ваш рік народження - ...|

\item \textbf{Ромб.} Дано довжини діагоналей ромба. Знайти його периметр.

{\itshape Намагайтеся супроводжувати введення та виведення текстовими повідомленнями.}

\item \textbf{Цифри числа.} Дано чотиризначне натуральне число. Переставити першу цифру до кінця числа. Вивести результат на екран.

\item \textbf{Формула.} Дані деякі дійсні числа $a$ та $b$. Обчислити значення виразу
$$
z=\cos(a^2+b^2)\cdot\dfrac{|a+b|}{2(b-2a)(a^2+b^2)}
$$
з точністю до сотих. Перевірка даних на коректність не потрібна.
\end{enumerate}


\cleardoublepage

\LabHeader{1}{ 5}

\begin{enumerate}
%\item Ввести с клавиатуры целое число (переменная \verb|n1|), вещественное число (переменная \verb|n2|) и слово (переменная \verb|w|).
%\begin{enumerate}
%\item Вывести значения переменных в одной строке, разделяя пробелом.
%\item Вывести значения переменных одно под другим.
%\item Вывести на экран значения переменных, сопровождая вывод комментариями, например:
%\begin{verbatim}
%Вы ввели три значения.
%Первое из них - ...
%Второе - ...
%А третье - ...
%\end{verbatim}
%\end{enumerate}

\item \textbf{Улюблений фільм.} Складіть програму, яка запитує у користувача назву його улюбленого фільму

(Приклад: \verb|Ваш улюблений фільм:|),

а потім після введення даних друкує на екрані повідомлення

(Приклад: \verb|Мій улюблений фільм - ...|).

\item \textbf{Пиріжок.} Пиріжок у місцевому буфеті коштує $a$ доларів $b$ центів. Скільки доларів і центів доведеться витратити на $n$ пиріжків?

{\itshape Намагайтеся супроводжувати введення та виведення текстовими повідомленнями.}

\item \textbf{Перестановка цифр.} Дане чотиризначне натуральне число. Дві його останні цифри було переставлено на початок числа. Вивести отримане число.

\item \textbf{Формула.} Наведено деякі дійсні числа $x$ і $y$. Обчислити значення виразу
$$
z=\dfrac{\sqrt{\mathstrut 1+\ln \left(1+x^2\right)}+\sqrt{\mathstrut 1+\dfrac{1}{1+|y|}}}{2x \,(1+y^2)}
$$
з точністю до сотих. Перевірка даних на коректність не потрібна.
\end{enumerate}

\cleardoublepage

\LabHeader{1}{ 6}

\begin{enumerate}
%\item Ввести с клавиатуры целое число (переменная \verb|n1|), вещественное число (переменная \verb|n2|) и слово (переменная \verb|w|).
%\begin{enumerate}
%\item Вывести значения переменных в одной строке, разделяя пробелом.
%\item Вывести значения переменных одно под другим.
%\item Вывести на экран значения переменных, сопровождая вывод комментариями, например:
%\begin{verbatim}
%Вы ввели три значения.
%Первое из них - ...
%Второе - ...
%А третье - ...
%\end{verbatim}
%\end{enumerate}

\item \textbf{Університет.} Складіть програму, яка запитує у користувача назву ВНЗ, в якій він навчається

(Приклад: \verb|У якому вузі Ви навчаєтеся?|),

а потім після введення даних друкує на екрані повідомлення

(Приклад: \verb|Я вчуся в ...|).

\item \textbf{Пробіжка.} Дано розміри (в метрах) футбольного поля, що має форму прямокутника. Яку відстань пробіжить спортсмен, що "намотав" кругів навколо цього поля? Число $n$ задається з клавіатури.

{\itshape Намагайтеся супроводжувати введення та виведення текстовими повідомленнями.}

\item \textbf{Задом наперед.} Дано тризначне натуральне число. Вивести на екран число, що виходить при прочитанні вихідного числа праворуч наліво.

\item \textbf{Формула.} Наведено деякі дійсні числа $a$ і $b$. Обчислити значення виразу
$$
z=\sin(a^2+b^2)\cdot\dfrac{\dfrac{a}{a^2+25}+b}{\sqrt{\mathstrut b}+\dfrac{a+b}{ 2}}
$$
з точністю до сотих. Перевірка даних на коректність не потрібна.
\end{enumerate}



\cleardoublepage



\LabHeader{1}{ 7}

\begin{enumerate}
%\item Ввести с клавиатуры целое число (переменная \verb|n1|), вещественное число (переменная \verb|n2|) и слово (переменная \verb|w|).
%\begin{enumerate}
%\item Вывести значения переменных в одной строке, разделяя пробелом.
%\item Вывести значения переменных одно под другим.
%\item Вывести на экран значения переменных, сопровождая вывод комментариями, например:
%\begin{verbatim}
%Вы ввели три значения.
%Первое из них - ...
%Второе - ...
%А третье - ...
%\end{verbatim}
%\end{enumerate}

\item \textbf{Привітання.} Складіть програму, яка запитує у користувача ім'я та прізвище

(Приклад: \verb|Як Вас звати?|),

а потім, після введення даних, друкує на екрані якесь привітання

(Приклад: \verb|Привіт, ... ...!|).

\item \textbf{Градуси та радіани.} Дано деяке число, що позначає градусну міру кута. Перевести це значення на радіани.

{\itshape Намагайтеся супроводжувати введення та виведення текстовими повідомленнями.}

\item \textbf{Перестановка цифр.} Дано тризначне натуральне число. Вивести на екран число, що виходить при перестановці другої та третьої цифр вихідного числа.

\item \textbf{Формула.} Дано деякі дійсні числа $x$ і $y$. Визначити значення віразу
$$
z=\ln \left(2-\sin x^2\right)\cdot\dfrac{x\,(1+x)}{|y|+\dfrac{1}{\sqrt{\mathstrut \mathstrut x ^2+10}}}
$$
з точністю до сотих. Перевірка даних на коректність не потрібна.
\end{enumerate}



\cleardoublepage



\LabHeader{1}{ 8}

\begin{enumerate}
%\item Ввести с клавиатуры целое число (переменная \verb|n1|), вещественное число (переменная \verb|n2|) и слово (переменная \verb|w|).
%\begin{enumerate}
%\item Вывести значения переменных в одной строке, разделяя пробелом.
%\item Вывести значения переменных одно под другим.
%\item Вывести на экран значения переменных, сопровождая вывод комментариями, например:
%\begin{verbatim}
%Вы ввели три значения.
%Первое из них - ...
%Второе - ...
%А третье - ...
%\end{verbatim}
%\end{enumerate}

\item \textbf{Улюблений герой.} Складіть програму, яка запитує користувача ім'я його улюбленого мультиплікаційного героя

(Приклад: \verb|Ваш улюблений мультяшний герой:|),

а потім після введення даних друкує на екрані повідомлення

(Приклад: \verb|Мій улюблений герой - ...|).

\item \textbf{Градуси Фаренгейта.} Відоме значення температури за шкалою Цельсія. Знайти відповідне
значення температури за шкалою Фаренгейта \ footnote {Для перерахунку за шкалою Фаренгейта необхідно вихідне значення температури помножити на ~1.8 і до результату додати 32.}.

{\itshape Намагайтеся супроводжувати введення та виведення текстовими повідомленнями.}

\item \textbf{Сума двох цифр.} Дано чотиризначне натуральне число. Знайти суму двох його останніх цифр.

\item \textbf{Формула.} Дані деякі дійсні числа $a$ та $b$. Обчислити значення виразу
$$
z=\cos(a^2+b^2)\cdot\dfrac{|a+b|}{2(b-2a)(a^2+b^2)}
$$
з точністю до сотих. Перевірка даних на коректність не потрібна.
\end{enumerate}

\cleardoublepage

\LabHeader{1}{ 9}

\begin{enumerate}
%\item Ввести с клавиатуры целое число (переменная \verb|n1|), вещественное число (переменная \verb|n2|) и слово (переменная \verb|w|).
%\begin{enumerate}
%\item Вывести значения переменных в одной строке, разделяя пробелом.
%\item Вывести значения переменных одно под другим.
%\item Вывести на экран значения переменных, сопровождая вывод комментариями, например:
%\begin{verbatim}
%Вы ввели три значения.
%Первое из них - ...
%Второе - ...
%А третье - ...
%\end{verbatim}
%\end{enumerate}

\item \textbf{Улюблений колір.} Складіть програму, яка запитує у користувача його улюблений колір

(Приклад: \verb|Назвіть Ваш улюблений колір:|),

а потім після введення даних друкує на екрані повідомлення

(Приклад: \verb|Мій улюблений колір - ...|).

\item \textbf{Щільність населення.} Відомі кількість жителів у місті та площа його території (у км$^2$). Визначити щільність населення в цьому місті Станом на липень 2018 року містом з найбільшою щільністю населення є Маніла "--- столиця держави Філіппіни (43\,079 осіб на км$^2$).}.

{\itshapeНамагайтесь супроводжувати введення та виведення текстовими повідомленнями.}

\item \textbf{Твір цифр.} Дано тризначне натуральне число. Знайти добуток його цифр.

\item \textbf{Формула.} Дані деякі дійсні числа $a$ та $b$. Обчислити значення виразу
$$
z=\cos(a^2+b^2)\cdot\dfrac{|a+b|}{2(b-2a)(a^2+b^2)}
$$
з точністю до сотих. Перевірка даних на коректність не потрібна.
\end{enumerate}

\cleardoublepage

\LabHeader{1}{ 10}

\begin{enumerate}
%\item Ввести с клавиатуры целое число (переменная \verb|n1|), вещественное число (переменная \verb|n2|) и слово (переменная \verb|w|).
%\begin{enumerate}
%\item Вывести значения переменных в одной строке, разделяя пробелом.
%\item Вывести значения переменных одно под другим.
%\item Вывести на экран значения переменных, сопровождая вывод комментариями, например:
%\begin{verbatim}
%Вы ввели три значения.
%Первое из них - ...
%Второе - ...
%А третье - ...
%\end{verbatim}
%\end{enumerate}

\item \textbf{Улюблений мультфільм.} Складіть програму, яка запитує користувача назва улюбленого мультфільму

(Приклад: \verb|Ваш улюблений мультфільм:|),

а потім після введення даних друкує на екрані повідомлення

(Приклад: \verb|Мій улюблений мультфільм - ...|).

\item \textbf{Квадрат і точки.} Яку найбільшу кількість точок з цілими координатами можна на листку в клітинку накрити квадратом зі стороною $n$ клітин? Число $n$ задається з клавіатури.

{\itshape Намагайтеся супроводжувати введення та виведення текстовими повідомленнями.}

\item \textbf{Номер квитка.} Дано шестизначне натуральне число "--- номер проїзного квитка. Знайти суму цифр у <<лівій>> та <<правій>> частинах номера.

\item \textbf{Формула.} Дані деякі дійсні числа $a$ та $b$. Обчислити значення виразу
$$
z=\sin(a^2-b^2)\cdot\dfrac{|a+b|}{2(b-2a)(a^2-b^2)}
$$
з точністю до сотих. Перевірка даних на коректність не потрібна.
\end{enumerate}

\cleardoublepage



\LabHeader{1}{ 11}

\begin{enumerate}
%\item Ввести с клавиатуры целое число (переменная \verb|n1|), вещественное число (переменная \verb|n2|) и слово (переменная \verb|w|).
%\begin{enumerate}
%\item Вывести значения переменных в одной строке, разделяя пробелом.
%\item Вывести значения переменных одно под другим.
%\item Вывести на экран значения переменных, сопровождая вывод комментариями, например:
%\begin{verbatim}
%Вы ввели три значения.
%Первое из них - ...
%Второе - ...
%А третье - ...
%\end{verbatim}
%\end{enumerate}

\item \textbf{Улюблений кіногерой.} Складіть програму, яка запитує користувача ім'я його улюбленого кіногероя

(Приклад: \verb|Ваш улюблений кіногерой:|),

а потім після введення даних друкує на екрані повідомлення

(Приклад: \verb|Мій улюблений кіногерой - ...|).

\item \textbf{Кути трикутника.} Дані катети прямокутного трикутника. Знайти усі його кути (у радіанах).
{\itshapeНамагайтеся супроводжувати введення та виведення текстовими повідомленнями.}

\item \textbf{Перестановка цифр.} Дано тризначне натуральне число. Вивести на екран число, що виходить при перестановці першої та другої цифр вихідного числа.

\item \textbf{Формула.} Дані деякі дійсні числа $a$ та $b$. Обчислити значення виразу
$$
z=\cos(a^2+b^2)\cdot\dfrac{|a+b|}{2(b-2a)(a^2+b^2)}
$$
з точністю до сотих. Перевірка даних на коректність не потрібна.
\end{enumerate}

\cleardoublepage

\LabHeader{1}{ 12}

\begin{enumerate}
%\item Ввести с клавиатуры целое число (переменная \verb|n1|), вещественное число (переменная \verb|n2|) и слово (переменная \verb|w|).
%\begin{enumerate}
%\item Вывести значения переменных в одной строке, разделяя пробелом.
%\item Вывести значения переменных одно под другим.
%\item Вывести на экран значения переменных, сопровождая вывод комментариями, например:
%\begin{verbatim}
%Вы ввели три значения.
%Первое из них - ...
%Второе - ...
%А третье - ...
%\end{verbatim}
%\end{enumerate}

\item \textbf{День тижня.} Складіть програму, яка запитує у користувача день тижня

(Приклад: \verb|Який сьогодні день тижня?|),

а потім після введення даних друкує на екрані повідомлення

(Приклад: \verb|Ура, сьогодні ...!|).

\item \textbf{Трикутник.} Дані катети прямокутного трикутника. Знайти його кути (у градусах).
{\itshape Намагайтеся супроводжувати введення та виведення текстовими повідомленнями.}

\item \textbf{Перестановка цифр.} Дане чотиризначне натуральне число. Поміняти місцями його першу та останню цифри. Вивести отримане число.

\item \textbf{Формула.} Дано деякі дійсні числа $x$ і $y$. Обчислити значення виразу
$$
z=\dfrac{\sqrt{\mathstrut 1+\ln \left(1+x^2\right)}+\sqrt{\mathstrut 1+\dfrac{1}{1+|y|}}}{2x \,(1+y^2)}
$$
з точністю до сотих. Перевірка даних на коректність не потрібна.
\end{enumerate}

\cleardoublepage

\LabHeader{1}{ 13}

\begin{enumerate}
%\item Ввести с клавиатуры целое число (переменная \verb|n1|), вещественное число (переменная \verb|n2|) и слово (переменная \verb|w|).
%\begin{enumerate}
%\item Вывести значения переменных в одной строке, разделяя пробелом.
%\item Вывести значения переменных одно под другим.
%\item Вывести на экран значения переменных, сопровождая вывод комментариями, например:
%\begin{verbatim}
%Вы ввели три значения.
%Первое из них - ...
%Второе - ...
%А третье - ...
%\end{verbatim}
%\end{enumerate}

\item \textbf{Улюблена група.} Складіть програму, яка запитує у користувача назву його улюбленої музичної групи

(Приклад: \verb|Ваша улюблена музична група:|),

а потім після введення даних друкує на екрані повідомлення

(Приклад: \verb|Моя улюблена група - ...|).

\item \textbf{Ранкова пробіжка.} Дано розміри (в метрах) футбольного поля, що має форму прямокутника. Скільки кіл потрібно << намотати >> навколо цього поля, щоб пробігти сумарно $ n $ км? Число $n$ задається з клавіатури.

{\itshape Намагайтеся супроводжувати введення та виведення текстовими повідомленнями.}

\item \textbf{Задом наперед.} Дано тризначне натуральне число. Вивести на екран число, що виходить при прочитанні вихідного числа праворуч наліво.

\item \textbf{Формула.} Дані деякі дійсні числа $a$ та $b$. Обчислити значення виразу
$$
z=\cos(a^2-b^2)\cdot\dfrac{\dfrac{a}{a^2+25}+b}{\sqrt{\mathstrut b}+\dfrac{a+b}{ 3}}
$$
з точністю до сотих. Перевірка даних на коректність не потрібна.
\end{enumerate}

\cleardoublepage

\LabHeader{1}{ 14}

\begin{enumerate}
%\item Ввести с клавиатуры целое число (переменная \verb|n1|), вещественное число (переменная \verb|n2|) и слово (переменная \verb|w|).
%\begin{enumerate}
%\item Вывести значения переменных в одной строке, разделяя пробелом.
%\item Вывести значения переменных одно под другим.
%\item Вывести на экран значения переменных, сопровождая вывод комментариями, например:
%\begin{verbatim}
%Вы ввели три значения.
%Первое из них - ...
%Второе - ...
%А третье - ...
%\end{verbatim}
%\end{enumerate}

\item \textbf{Улюблене місто.} Складіть програму, яка запитує у користувача назву улюбленого міста

(Приклад: \verb|Яке Ваше улюблене місто?|),

а потім, після введення даних, друкує на екрані якесь привітання

(Приклад: \verb|Моє улюблене місто - ...|).

\item \textbf{Біг по колу.} Леді Гамільтон "--- дуже спортивна жінка, незважаючи на поважний вік. Її улюблене заняття зранку"--- пробіжка навколо знаменитого фонтану в парку Кентербері разом з вірним псом на прізвисько Лінкольн. Яка відстань (в метрах) зазвичай пробігає вранці місіс Гамільтон з Лінкольном, якщо радіус фонтану $r$ метрів і вони роблять $n$ кіл?

{\itshape Намагайтеся супроводжувати введення та виведення текстовими повідомленнями.}

\item \textbf{Сума квадратів цифр.} Дано тризначне натуральне число. Вивести на екран суму квадратів його цифр.

\item \textbf{Формула.} Дано деякі дійсні числа $x$ і $y$. Обчислити значення виразу
$$
z=\ln \left(2-\sin x^2\right)\cdot\dfrac{x\,(1+x)}{|y|+\dfrac{1}{\sqrt{\mathstrut \mathstrut x ^2+10}}}
$$
з точністю до сотих. Перевірка даних на коректність не потрібна. 
\end{enumerate}

\cleardoublepage

\LabHeader{1}{ 15}

\begin{enumerate}
%\item Ввести с клавиатуры целое число (переменная \verb|n1|), вещественное число (переменная \verb|n2|) и слово (переменная \verb|w|).
%\begin{enumerate}
%\item Вывести значения переменных в одной строке, разделяя пробелом.
%\item Вывести значения переменных одно под другим.
%\item Вывести на экран значения переменных, сопровождая вывод комментариями, например:
%\begin{verbatim}
%Вы ввели три значения.
%Первое из них - ...
%Второе - ...
%А третье - ...
%\end{verbatim}
%\end{enumerate}

\item \textbf{Привітання.} Складіть програму, яка запитує у користувача його ім'я та рік народження, а потім виводить повідомлення

\verb|Здрастуйте, ...! Ваш рік народження - ...|

\item \textbf{Брутальна клумба.} Знайти площу клумби, якщо вона є ідеальним колом, а на її огородження пішло $n$ метрів бетонного паркану. Висоту паркану не враховуємо. Кількість $n$ вводиться з клавіатури.

{\itshape Намагайтеся супроводжувати введення та виведення текстовими повідомленнями.}

\item \textbf{Цифри числа.} Дано чотиризначне натуральне число. Переставити останню цифру на початку числа. Вивести результат на екран.

\item \textbf{Формула.} Дані деякі дійсні числа $a$ та $b$. Обчислити значення виразу
$$
z=\cos(a^2+b^2)\cdot\dfrac{|a+b|}{2(b-2a)(a^2+b^2)}
$$
з точністю до сотих. Перевірка даних на коректність не потрібна.
\end{enumerate}

\cleardoublepage

\LabHeader{1}{ 16}

\begin{enumerate}
%\item Ввести с клавиатуры целое число (переменная \verb|n1|), вещественное число (переменная \verb|n2|) и слово (переменная \verb|w|).
%\begin{enumerate}
%\item Вывести значения переменных в одной строке, разделяя пробелом.
%\item Вывести значения переменных одно под другим.
%\item Вывести на экран значения переменных, сопровождая вывод комментариями, например:
%\begin{verbatim}
%Вы ввели три значения.
%Первое из них - ...
%Второе - ...
%А третье - ...
%\end{verbatim}
%\end{enumerate}

\item \textbf{Улюблений автор.} Складіть програму, яка запитує у~користувача прізвище його улюбленого автора книг

(Приклад: \verb|Ваш улюблений автор:|),

а потім після введення даних друкує на екрані повідомлення

(Приклад: \verb|Мій улюблений автор - ...|).

\item \textbf{Сума чисел.} Знайти суму всіх цілих чисел на проміжку $[a; b] $. Числа $a$ і $b$ "--- цілі, що вводяться з клавіатури.

{\itshape Намагайтеся супроводжувати введення та виведення текстовими повідомленнями.}

\item \textbf{Сума двох цифр.} Дано чотиризначне натуральне число. Знайти суму двох його останніх цифр.

\item \textbf{Формула.} Дані деякі дійсні числа $a$ та $b$. Обчислити значення виразу
$$
z=\sin(a^2+b^2)\cdot\dfrac{\dfrac{a}{a^2+25}+b}{\sqrt{\mathstrut b}+\dfrac{a+b}{ 2}}
$$
з точністю до сотих. Перевірка даних на коректність не потрібна. 
\end{enumerate}

\cleardoublepage

\LabHeader{1}{ 17}

\begin{enumerate}
%\item Ввести с клавиатуры целое число (переменная \verb|n1|), вещественное число (переменная \verb|n2|) и слово (переменная \verb|w|).
%\begin{enumerate}
%\item Вывести значения переменных в одной строке, разделяя пробелом.
%\item Вывести значения переменных одно под другим.
%\item Вывести на экран значения переменных, сопровождая вывод комментариями, например:
%\begin{verbatim}
%Вы ввели три значения.
%Первое из них - ...
%Второе - ...
%А третье - ...
%\end{verbatim}
%\end{enumerate}

\item \textbf{Улюблене місто.} Складіть програму, яка запитує у користувача назву улюбленого міста

(Приклад: \verb|Яке Ваше улюблене місто?|),

а потім, після введення даних, друкує на екрані якесь привітання

(Приклад: \verb|Моє улюблене місто - ...|).

\item \textbf{Куб.} Дана довжина ребра куба. Знайти об'єм куба та площу його бічної поверхні.

{\itshape Намагайтеся супроводжувати введення та виведення текстовими повідомленнями.}

\item \textbf{Перестановка цифр.} Дано чотиризначне натуральне число. Поміняти місцями його першу та останню цифри. Вивести отримане число.

\item \textbf{Формула.} Дані деякі дійсні числа $a$ та $b$. Обчислити значення виразу
$$
z=\cos(a^2+b^2)\cdot\dfrac{|a+b|}{2(b-2a)(a^2+b^2)}
$$
з точністю до сотих. Перевірка даних на коректність не потрібна.  
\end{enumerate}
